% =====================================================================
%  MARCO REFERENCIAL (versión recortada, para comparar con
%  capitulos/06-marco-referencial.tex)
%
%  Recorte del 2026-10-24, para llevar el documento a 60-80 páginas:
%   - Antecedentes: descripciones más cortas; la distribución del cómputo
%     pasa a la subsección del aporte, con sus tres citas.
%   - Marco conceptual: sin subsecciones. Sigue en párrafos corridos, como
%     pidió el director el 2026-10-05, uno por cada categoría del
%     anteproyecto; los términos que el marco teórico desarrolla se definen
%     en una línea y remiten a él.
%   - Marco teórico: salen los hechos del proyecto que ya cuenta el
%     desarrollo (carga de la tarjeta, las dos distribuciones de ROS 2, la
%     banda muerta, el caso del criterio de continuidad, el intervalo de la
%     campaña). El seguimiento de la trayectoria se une a la planificación.
%   - Marco legal: sin cambios.
%   - Marco ambiental: dos párrafos, sin subsecciones.
%  Ninguna cita sale del capítulo.
%
%  REGISTRO: lectura de pregrado. Frases cortas, voz activa, una idea por
%  frase. Los términos van en negrita solo en el marco conceptual, donde
%  sirven para barrer la sección con la vista (pedido del director).
% =====================================================================
\chapter{Marco referencial}
\label{cap:marco}

% =====================================================================
\section{Antecedentes}
\label{sec:antecedentes}
% =====================================================================

\subsection{Robots que guían a personas}

\textcite{burgard1999experiences} describen RHINO, un robot guía que operó de
forma autónoma durante seis días en un museo en 1997. Integraba localización,
construcción de mapas, evasión de obstáculos, planificación, interacción con los
visitantes y una interfaz web para seguirlo a distancia. De ese trabajo se toma
una idea vigente: guiar personas es un problema de integración tanto como de
navegación, porque ninguna pieza sirve si las demás fallan. Del mismo grupo salió
la localización de Monte Carlo \autocite{fox1999monte}, base del algoritmo de
localización de los agentes de este proyecto.

\textcite{guerreiro2019cabot} llevan la idea a la asistencia de personas ciegas.
Su robot, CaBot, tiene forma de maleta y transmite la dirección por un mango que
vibra; lo evaluaron con diez participantes ciegos. El trabajo muestra dos
exigencias que también aparecen aquí: el robot tiene que adaptarse a una persona
que camina a su lado, y la interfaz tiene que comunicar el estado del guiado sin
exigir atención constante.

\textcite{li2023bug} abordan la seguridad de compartir el espacio con personas,
modificando el algoritmo Bug con sensores de proximidad para reducir las
colisiones con los trabajadores cercanos. En este proyecto la respuesta está en
la operación: una persona acompaña a cada agente durante las maniobras, y el
usuario puede cancelar la misión en cualquier momento.
\textcite{jimenez2024arquitectura} integra al usuario en el propio sistema de
mapas, con representaciones semánticas y reconocimiento de personas, en una
plataforma asistencial de varios niveles. Es el antecedente más próximo en
intención, porque trata a la persona como parte del problema de navegación.

En los tres casos un mismo robot acompaña a la persona durante todo el recorrido.
Ese supuesto es el que este trabajo cambia.

\subsection{Navegación autónoma en interiores}

\textcite{sojan2025autonomous} integran percepción, localización y planificación
sobre ROS~2 para desplazarse en interiores dinámicos, con un mapa plano y
localización por filtro de partículas. Es la arquitectura que este proyecto toma
como punto de partida en cada agente.

El sistema que implementa ese flujo en ROS~2 es Nav2.
\textcite{macenski2020marathon} lo presentan como sucesor de la pila de
navegación de la primera generación de ROS: organiza la tarea con un árbol de
comportamiento y arranca cada servidor como un nodo de ciclo de vida. Lo
validaron en un campus universitario, navegando entre estudiantes.
\textcite{macenski2023survey} comparan los algoritmos de planificación, control y
localización que hoy ofrece el ecosistema. Las dos referencias sirvieron para elegir los componentes de navegación.
\textcite{varghese2023multirobot} llevan el problema a varios agentes que
construyen un mapa compartido y mantienen consistente su localización, con la
misma técnica bidimensional sobre ROS.

Otra línea resuelve la navegación con aprendizaje en lugar de planificación
explícita: redes neuronales que aprenden a conducir solo con imágenes
\autocite{bojarski2016end}, a evadir obstáculos fusionando sensores
\autocite{zain2025obstacle} o a reaccionar combinando aprendizaje profundo y lógica
difusa \autocite{lopez2024enhancing}. Este proyecto no sigue esa vía, porque
exige datos etiquetados del edificio, que no existen, y porque con planificación
clásica cada decisión del agente se puede inspeccionar cuando hay que explicar
un fallo.

\subsection{Coordinación entre varios robots}

Las estrategias de coordinación revisadas se diferencian por el nivel en el que
cooperan. \textcite{nair2024collaborative} comparten la información sensorial
entre los robots para mejorar el desempeño del conjunto.
\textcite{roszkowska2023multilevel} proponen una estructura jerárquica que
supervisa el desplazamiento simultáneo de varios agentes para prevenir
conflictos. \textcite{rudrusamy2022broadcast} resuelven el problema desde el
canal de comunicación, con control por difusión y seguimiento visual.
\textcite{xu2025modular} aportan una arquitectura modular sobre ROS para
asignación cooperativa de tareas, y \textcite{zhan2025kubernetes} estudian la
resiliencia del sistema cuando un nodo falla.

El marco teórico de la asignación lo fijan \textcite{gerkey2004taxonomy}, cuya
taxonomía clasifica cualquier esquema de reparto según cuántas tareas acepta un
robot, cuántos robots requiere una tarea y con cuánta anticipación se decide. Esa taxonomía, con su extensión por \textcite{korsah2013comprehensive}, es la que sitúa el esquema de asignación de este trabajo.

El antecedente más cercano es \textcite{srivastava2025deliver}, que coordinan
recogida y entrega entre varios robots mediante planificación con relevos. Allí
lo que pasa de un agente a otro es un paquete; aquí es una persona. Un paquete
puede esperar indefinidamente en el punto de transferencia, y una persona no. Por
eso el relevo de este trabajo incluye una confirmación explícita del usuario al
llegar al otro piso, con un aviso a los 60~s y un plazo máximo de 120~s. De
alcance comparable, por ser también un proyecto de grado de ingeniería, es
\textcite{piguave2023multirobot}, un sistema multirrobot de código abierto sobre
ROS~2. Según la revisión del anteproyecto, todos estos trabajos operan en un solo
nivel físico.

\subsection{Movilidad entre pisos}

El cambio de nivel se ha abordado sobre todo desde la logística y la
exploración. \textcite{kim2024indoor} combinan un mapa por piso con un grafo
topológico que los conecta, e integran la interacción con el ascensor.
\textcite{song2024elevator} automatizan el uso del elevador con dispositivos
conectados instalados en el edificio. \textcite{luo2025elevator} deciden entre
ascensor y escalera con aprendizaje por refuerzo jerárquico, según el tránsito
que encuentra el robot. \textcite{chen2025lite} descomponen el entorno en pisos y
escaleras para extender las técnicas bidimensionales a varios niveles, y
\textcite{gong2026stairway} incorporan conciencia de piso a la navegación hacia
objetos sin entrenamiento previo del escenario.

Estas propuestas atacan la movilidad entre pisos desde la infraestructura, el
modelado topológico y la decisión adaptativa, y comparten un supuesto: el robot
cruza de nivel, subiendo o accionando el ascensor. Todas se orientan a entrega o
exploración antes que a la asistencia de personas.

\subsection{Aporte de este trabajo frente a los antecedentes}
\label{sec:aporte}

A esas líneas se suma la de dónde ejecutar el procesamiento.
\textcite{kehoe2015cloud} revisan las arquitecturas que lo delegan a servidores
remotos, \textcite{tahir2025edge} las que lo acercan al borde de la red, y
\textcite{malavolta2023offloading} miden el costo de esa delegación sobre WiFi en
robots terrestres. Sus resultados llevaron a que la navegación corra completa en cada agente y solo la coordinación, que intercambia mensajes pequeños,
se centralice en un computador portátil, sin depender de salida a internet.

La literatura cubre bien cada componente por separado. La
tabla~\ref{tab:antecedentes} resume lo que cada línea aporta a este proyecto y el
supuesto que este proyecto no comparte.

\begin{table}[H]
  \centering\small
  \caption{Líneas de antecedentes, lo que aportan a este trabajo y el supuesto
    que este trabajo cambia.}
  \label{tab:antecedentes}
  \begin{tabular}{@{}p{0.2\textwidth}p{0.36\textwidth}p{0.36\textwidth}@{}}
    \toprule
    \textbf{Línea} & \textbf{Lo que se toma de ella} & \textbf{Supuesto que no se comparte} \\ \midrule
    Robots guía & La integración de navegación e interacción; la localización de Monte Carlo & Un solo robot acompaña todo el recorrido \\ \addlinespace
    Navegación en interiores & El flujo mapa, localización y planificación sobre ROS~2 (Nav2) & Ninguno: se adopta, adaptado a la cinemática del agente \\ \addlinespace
    Coordinación multirrobot & La taxonomía de asignación; la planificación con relevos & Lo que se traspasa puede esperar sin límite \\ \addlinespace
    Movilidad entre pisos & El mapa por piso unido por un punto de transición & El robot cruza de nivel \\ \addlinespace
    Distribución del cómputo & La navegación a bordo y la coordinación en un nodo central & --- \\ \bottomrule
  \end{tabular}
\end{table}

Lo que no aparece en ninguna línea es la combinación que aquí se propone:
resolver la discontinuidad entre pisos sin que ningún agente la atraviese. El
recorrido se reparte entre dos agentes dedicados, cada uno a un nivel, y la
persona pasa de uno a otro en la zona de transición. La diferencia con los
trabajos de movilidad vertical es económica antes que algorítmica. Esos trabajos
obtienen la continuidad con hardware, con mecanismos para escaleras o con
intervención del ascensor. Este proyecto la obtiene con comunicación entre
agentes, que es lo que los trabajos de coordinación ya resuelven, aplicada a un
problema donde no se había usado.

% =====================================================================
\section{Marco conceptual}
\label{sec:conceptual}
% =====================================================================
% FORMA: parrafos corridos, no listas de definiciones (pedido del director,
% 2026-10-05). Un parrafo por cada categoria del anteproyecto. Los terminos
% siguen en negrita para poder barrer la seccion con la vista.
% =====================================================================

En navegación y percepción, la \textbf{navegación autónoma} es el desplazamiento
de un robot sin intervención humana, y reúne tres capacidades: percibir el
entorno, estimar dónde está y decidir hacia dónde moverse. Cada agente percibe
con un \textbf{LiDAR}, un sensor láser que mide distancias en un plano, en
360$^\circ$ a su alrededor, y con una \textbf{unidad de medición inercial} o
IMU, que mide la velocidad de giro y la aceleración. Como no tiene codificadores
en las ruedas, su \textbf{odometría} sale de comparar barridos sucesivos del
láser, y \textbf{AMCL} la corrige contra un mapa conocido. Que la odometría
tenga con qué estimar el avance depende del sitio. Para medirlo, el proyecto definió la \textbf{información de avance} de un
barrido: la fracción de los rayos del LiDAR que inciden sobre superficies
orientadas hacia la dirección de marcha. Solo esos rayos cambian de longitud
cuando el agente avanza. El entorno
se representa con un \textbf{mapa de ocupación}, que dice qué celdas están libres,
y con \textbf{mapas de costos}, que dan a cada celda un costo de tránsito. La
\textbf{zona de transición} es el lugar del edificio donde se cambia de nivel; en
este sistema, la escalera.

En control y locomoción, el agente sigue un \textbf{modelo cinemático tipo
automóvil}, o de \textbf{dirección Ackermann}: dirige con el eje delantero y
tiene un \textbf{radio de giro mínimo}. No puede girar en el sitio, como sí lo
hace un robot de tracción diferencial. Para invertir el rumbo en un
pasillo angosto hace una \textbf{media vuelta en varios tiempos}, alternando
marcha atrás y avance con la dirección a tope. Cada punto de una trayectoria en
el que el agente invierte la marcha es una \textbf{cúspide}. Las órdenes llegan
al motor y al servomotor de dirección como señales de \textbf{modulación por
ancho de pulso} (PWM). La tracción funciona en lazo abierto, sin un
\textbf{control PID} que regule la velocidad, y tiene una \textbf{banda muerta}:
por debajo de 0,40~m/s el motor no mueve al agente. Al final de cada tramo, el
\textbf{verificador de meta} decide si el agente llegó, comparando su posición
con una tolerancia.

En comunicación y sistemas distribuidos, el \textbf{middleware robótico} del
proyecto es \textbf{ROS~2}. Organiza el software en \textbf{nodos}, programas que
cumplen cada uno una función, que se comunican por \textbf{tópicos} de
publicación y suscripción, \textbf{servicios} de petición y respuesta y
\textbf{acciones}. Una acción representa una tarea larga: tiene una meta,
retroalimentación periódica y un resultado, y se puede cancelar. El nodo que la
ofrece es el \textbf{servidor de acción} y el que la invoca, el
\textbf{cliente}; el coordinador de este proyecto es cliente ante cada agente y
servidor ante la interfaz. Por debajo, ROS~2 usa \textbf{DDS}, un estándar en el
que los participantes de la red se encuentran solos por \textbf{descubrimiento}.
Quién habla con quién lo deciden el \textbf{dominio}, un número que separa el
descubrimiento; la \textbf{partición}, un filtro dentro de un dominio, y el
\textbf{espacio de nombres}, un prefijo como \texttt{/robot1} que indica a qué
agente pertenece cada tópico. Las políticas de entrega de cada canal se llaman
\textbf{calidad de servicio}. El \textbf{puente \textit{rosbridge}} traduce entre
ROS~2 y mensajes JSON sobre WebSocket, y un \textbf{contenedor} es un entorno de
ejecución aislado, con sus propias bibliotecas, que permite tener dos versiones
de ROS~2 en un mismo computador.

En arquitectura del sistema, la \textbf{asignación de tareas} es el reparto de
actividades entre los agentes según su disponibilidad y su ubicación. Cuando una
tarea cruza de un nivel a otro, el reparto toma la forma de un
\textbf{relevo}: un agente entrega la responsabilidad de guiar al usuario y el
otro la recibe en el \textbf{punto de transferencia}. La \textbf{máquina de
estados de misión} describe en qué punto va cada misión, y sus estados se
publican para la interfaz y para el registro. El \textbf{planificador puro}
calcula los tramos de una misión a partir del catálogo de puntos y del par
origen-destino. No depende del middleware ni de ningún estado externo, y por eso
se puede verificar con cientos de casos sin levantar el simulador. El
\textbf{registro de misión} es un documento estructurado que se compone a partir
de la grabación de cada ejecución y se valida contra un esquema; es la evidencia
primaria del experimento.

En interacción humano-robot, la \textbf{interacción humano-robot} o HRI estudia
el diseño y la evaluación de sistemas robóticos pensados para trabajar con
personas \autocite{jimenez2024arquitectura}. Aquí toma la forma de una
\textbf{interfaz gráfica} web desde la que el usuario hace la \textbf{solicitud
de servicio}. La \textbf{interfaz responsiva de selección} permite buscar un
destino escribiendo su nombre sin perder el catálogo completo, y la
\textbf{retroalimentación literal} es el principio por el que la interfaz
muestra tal cual el texto que redacta el coordinador.

En métricas y criterios de evaluación, las cuatro métricas son las del
anteproyecto: el \textbf{tiempo de respuesta}, el \textbf{tiempo de asignación}, la \textbf{tasa de éxito} y la
\textbf{continuidad de servicio}. El tiempo de asignación es el intervalo entre
aceptar la solicitud y designar el agente que la atiende. Para medirlas se usan
además el \textbf{intervalo de confianza de Wilson} y la \textbf{falsabilidad}
de un criterio. Un \textbf{invariante estructural} es una
propiedad que se cumple siempre por el diseño del sistema y que, por eso, no es
un resultado experimental. La \textbf{cota superior por misión} es lo que se
informa cuando el instrumento no tiene resolución para dar el valor exacto. Con los agentes físicos, la demostración se organizó en \textbf{pruebas de
verificación}, condiciones fijadas antes de correrlas, como
que la odometría registre el avance con un error del 10~\% o menos. Decidirlas
exige una \textbf{verdad de terreno}, la medida independiente con que se
contrasta el sistema, que en este trabajo es la distancia tomada con flexómetro.

% =====================================================================
\section{Marco teórico}
\label{sec:teorico}
% =====================================================================

\subsection{Asignación de tareas en sistemas multirrobot}
\label{sec:teo-asignacion}

\textcite{gerkey2004taxonomy} clasifican los problemas de asignación de tareas
entre robots según tres ejes. El primero distingue robots que ejecutan una sola
tarea a la vez (ST, \textit{single-task}) de robots que atienden varias (MT). El
segundo distingue tareas que requieren un solo robot (SR, \textit{single-robot})
de tareas que requieren varios (MR). El tercero distingue la asignación
instantánea (IA), que se decide al llegar cada tarea, de la extendida en el
tiempo (TA), que planifica asignaciones futuras.

La asignación de este sistema es ST--SR--IA. Cada agente ejecuta un tramo a la vez, cada tramo lo ejecuta uno solo, y la elección se hace al llegar la solicitud. Los mismos autores demuestran que esa clase es una instancia del
problema de asignación óptima, que el método húngaro resuelve en tiempo
polinómico \autocite{kuhn1955hungarian}. En este problema, sin embargo, la matriz
de costos tiene un solo candidato admisible por tarea, porque cada agente está
confinado a su nivel y hay uno solo por nivel. El método húngaro devuelve siempre
la correspondencia que el sistema guarda como tabla de configuración, en la que
cada nivel tiene asignado su agente. La ausencia de un optimizador es, por
tanto, una consecuencia de la estructura del problema y no una carencia de la
implementación.

Por la misma razón se descartaron los mecanismos de mercado. En el protocolo de
red de contratos \autocite{smith1980contract}, un gestor anuncia la tarea, los
agentes ofertan su costo y la tarea se adjudica a la mejor oferta. MURDOCH
\autocite{gerkey2002sold} aplica ese esquema a robots sobre publicación y
suscripción, y CBBA \autocite{choi2009consensus} lo lleva a una subasta
descentralizada en la que los agentes llegan a consenso sin coordinador central.
Con un solo agente elegible por tramo, cualquier subasta tiene un único postor:
añade mensajes y latencia sin cambiar la decisión. CBBA, además, elimina el
coordinador central, que la arquitectura mantiene por diseño.

La misión completa sí se sale de esa clase. El relevo impone una restricción de
precedencia: el tramo del segundo agente no puede empezar hasta que el primero
termine el suyo y el usuario cambie de piso. Las agendas de los dos agentes
quedan acopladas, y la taxonomía de 2004 está construida para tareas
independientes. \textcite{korsah2013comprehensive} la extienden según el grado
de interdependencia entre tareas, y en su clasificación esta misión pertenece a
la clase XD, de dependencias entre agendas (\textit{cross-schedule
dependencies}). Ahí está el aporte del trabajo: la asignación es trivial y óptima
por construcción, y la coordinación del relevo pertenece a una clase de problema
que la taxonomía clásica del área no contempla.

El análisis también delimita el alcance. El sistema atiende una misión a la vez y
no tiene noción de agente ocupado ni cola de solicitudes. Con dos o más agentes
elegibles por tarea, o con misiones simultáneas, la matriz de costos deja de
tener un único candidato, y la elección del agente pasa a depender del método de
asignación. Esa
extensión se plantea en los trabajos futuros.

\subsection{ROS~2 y la comunicación entre los componentes}
\label{sec:teo-ros2}

Todos los componentes se comunican mediante ROS~2, que
\textcite{macenski2022ros2} describen como un rediseño sobre DDS, el estándar de
comunicación de datos en tiempo real del Object Management Group
\autocite{omg2015dds}. DDS aporta el descubrimiento automático de los
participantes y las políticas de calidad de servicio de cada canal.

El descubrimiento tiene un costo que crece con el tamaño del sistema. Por
omisión, cada proceso se anuncia en la red y empareja sus publicadores y
suscriptores con los de todos los demás, y sobre una red inalámbrica ese tráfico
compite con la navegación por el procesador de cada agente. ROS~2 Jazzy permite restringir el
descubrimiento a los procesos del mismo computador y declarar una lista fija de
equipos con los que sí se comunica. Con esas dos opciones la red adopta la
topología en estrella del diseño. Dentro de cada agente los procesos se
descubren entre sí, y hacia afuera solo se anuncia al computador portátil lo que
el coordinador necesita, de modo que los agentes no se ven entre sí.

Dos mecanismos más separan a los agentes. La partición de comunicaciones deja
dentro de cada agente los tópicos del software del fabricante, que no llevan
espacio de nombres. Un espacio de nombres distinto para cada agente distingue la navegación, el estado y la media vuelta de cada uno cuando llegan al portátil. El sistema usa además dos distribuciones de ROS~2, Humble en simulación y Jazzy en los agentes físicos.

\subsection{Cinemática tipo automóvil, planificación y seguimiento de trayectorias}
\label{sec:teo-planificacion}

Un vehículo con dirección Ackermann está sujeto a una restricción no holonómica:
no puede desplazarse de lado ni girar sobre su propio eje, y su curvatura está
acotada por el radio de giro mínimo \autocite{siegwart2011introduction}. Por eso
no toda trayectoria libre de obstáculos es ejecutable, y un planificador que solo
considera la posición propone giros que el vehículo no puede hacer.

\textcite{reeds1990optimal} resolvieron el problema geométrico de fondo: el
camino más corto entre dos poses para un automóvil que puede avanzar y
retroceder, con radio de giro mínimo. Basta considerar caminos formados por arcos
de radio mínimo y segmentos rectos, con a lo sumo dos inversiones de marcha, y el
óptimo está siempre en un conjunto finito de candidatos con fórmulas explícitas.
Como los agentes de este trabajo pueden retroceder, ese es su modelo de movimiento.

Sobre un mapa con obstáculos hace falta además una búsqueda.
\textcite{dolgov2010path} propusieron el A* híbrido: una búsqueda A* sobre la
posición y la orientación del vehículo, en la que cada expansión respeta la
cinemática y el camino resultante se suaviza después. Lo validaron en los
estacionamientos del DARPA Urban Challenge de 2007, un escenario de maniobras con
inversión de marcha parecido al de un pasillo angosto. Nav2 lo implementa en el
planificador Smac, que admite el modelo Reeds-Shepp \autocite{macenski2024smac},
y es el que usan los agentes de este proyecto. El planificador trabaja sobre
mapas de costos organizados por capas \autocite{lu2014layered}: la estática viene
del mapa, la de obstáculos, de las lecturas del sensor, y la de inflado extiende
el costo alrededor de cada obstáculo según el tamaño del agente. El mapa de
costos global sirve al planificador y el local, que cubre los alrededores del
agente, al controlador.

El controlador convierte la trayectoria en órdenes de velocidad. El algoritmo Pure Pursuit, o de persecución pura \autocite{coulter1992implementation}, elige sobre la trayectoria un punto
situado a una distancia de anticipación $\ell$ del agente y calcula el arco de
circunferencia que lleva hasta él. Llamando $y$ a la distancia lateral del punto
al eje del agente, la curvatura ordenada es
\begin{equation}
  \kappa = \frac{2\,y}{\ell^{2}} .
\end{equation}
No requiere un modelo dinámico y siempre ordena arcos que la dirección puede
seguir. \textcite{macenski2023rpp} le añadieron una regulación de velocidad en
las curvas cerradas y cerca de los obstáculos, y la marcha atrás, de modo que sigue las cúspides que propone el planificador; esa versión, Regulated Pure Pursuit, es la que usa este trabajo.

Esa configuración tiene un límite cerca de la meta. En este trabajo la distancia
de anticipación va de 0,3 a 0,9~m. Cuando una cúspide queda muy cerca de la
meta, esa distancia se acorta por debajo de un umbral en el que una guarda de seguridad del
propio controlador anula la curvatura. El agente sigue invirtiendo la marcha
sin dirigir y maniobra decenas de veces en el mismo punto. Por eso la media
vuelta se resolvió con una maniobra propia del agente.

\subsection{Localización, odometría y observabilidad}
\label{sec:teo-localizacion}

La localización de Monte Carlo representa la creencia sobre la pose del agente
con un conjunto de partículas, cada una con un peso
\autocite{fox1999monte,thrun2005probabilistic}. En cada paso, el modelo de
movimiento desplaza las partículas según la odometría. Después, el modelo de
medición las pondera según cuánto se parece el barrido del láser a lo que se vería desde esa
pose en el mapa, y un remuestreo las concentra en las poses más probables. AMCL
es la variante que ajusta el número de partículas a la incertidumbre; en este
trabajo, entre 500 y 2000. Su precisión depende de la calidad de la odometría
entre una actualización y la siguiente.

Como los agentes no tienen codificadores en las ruedas, la odometría sale del
láser. \textcite{jaimez2016planar} estiman el movimiento plano de un láser a
partir de barridos consecutivos con un enfoque de flujo de distancias
(\textit{range flow}). Alinean los barridos de forma densa, usando sus
gradientes y sin buscar correspondencias entre puntos, en alrededor de un
milisegundo por barrido. Es el método que implementa el paquete rf2o, que usan
los agentes.

Todo estimador que se apoya en la geometría del entorno hereda una limitación.
\textcite{zhang2016degeneracy} la formalizan como degeneración: hay direcciones
del espacio de estados en las que las mediciones no restringen la solución. Un
pasillo recto y de sección uniforme es el caso típico. Al avanzar a lo largo de
su eje, cada barrido es casi idéntico al anterior y el estimador devuelve un
desplazamiento cercano a cero, coherente con lo que mide el sensor. La posición
longitudinal es inobservable.

De ahí sale la diferencia entre dos tipos de error que se parecen. Un error de
calibración se corrige ajustando parámetros; una inobservabilidad estructural no,
porque la información no está en la señal. La información de avance mide cuánta hay en cada sitio. El criterio que
el proyecto adoptó para elegir dónde operar es que el entorno ofrezca estructura
no uniforme dentro del alcance del sensor durante todo el recorrido. La
construcción de mapas comparte el límite: la herramienta de SLAM de ROS~2
\autocite{macenski2021slam} usa la misma odometría del láser, y en un pasillo
uniforme no tiene con qué cerrar el mapa. Por eso los mapas de este proyecto se
derivaron de un levantamiento con flexómetro.

El rumbo tiene una solución distinta, porque el giroscopio de la IMU lo mide de
forma directa, sin depender del entorno. Un filtro de Kalman extendido combina
las dos fuentes, ponderando cada medición según su incertidumbre
\autocite{thrun2005probabilistic}. El paquete robot\_localization lo implementa
para ROS de forma general, permitiendo elegir qué componente de cada sensor se
fusiona \autocite{moore2014generalized}. En este trabajo el filtro toma de la
odometría láser las velocidades de avance y de la IMU la velocidad de giro. El
rumbo mejora; la posición no, porque sigue saliendo del láser.

\subsection{Simulación del sistema}
\label{sec:teo-simulacion}

La campaña de evaluación corrió en Gazebo~11 con ROS~2 Humble. Gazebo es un
simulador de código abierto pensado desde su origen para varios robots, con motor
de física, modelos de sensores e integración con ROS \autocite{koenig2004gazebo}.
La cinemática Ackermann del agente se simula con un complemento propio del
paquete de simulación del DeepRacer, en lugar del genérico de Gazebo. Los dos agentes simulados comparten un dominio de comunicaciones y se separan por espacio de
nombres, incluido el marco de referencia de su mapa. Con la configuración por
omisión, dar a cada agente un dominio distinto corta también su descubrimiento con el coordinador. El tiempo que ven los nodos avanza en saltos
de 100~ms, y eso fija la resolución con que se miden los tiempos de la campaña.

\subsection{La interfaz web como cliente del sistema}
\label{sec:teo-interfaz}

\textcite{crick2017rosbridge} propusieron \textit{rosbridge} para que programas
ajenos a ROS, como un navegador, usen tópicos, servicios y acciones mediante
mensajes JSON sobre WebSocket. La versión para ROS~2 expone las acciones con
operaciones propias, que la biblioteca de cliente web más difundida no
implementa, porque fue escrita para las acciones de la primera generación de ROS.
Por eso la interfaz de este trabajo es un cliente propio, sin dependencias
externas, ya que debe funcionar sin salida a internet.

La interfaz no contiene lógica del dominio: el texto que lee el usuario en cada
estado lo escribe el coordinador, y la interfaz lo muestra literal. Si el
navegador interpretara el estado para decidir qué texto mostrar, duplicaría una
lógica que ya vive en el planificador, y las dos copias acabarían divergiendo.

\subsection{Evaluación experimental}
\label{sec:teo-evaluacion}

La evaluación se apoya en las cuatro métricas del anteproyecto, cada una con una
definición operativa que indica qué marca del registro se resta de cuál. El
protocolo se escribió antes de instrumentar el sistema, y las misiones se sortean
con una semilla fijada antes de correrlas, para que ni el instrumento ni los
recorridos se elijan después de ver los resultados.

La tasa de éxito es una proporción estimada sobre una muestra pequeña, y para
ella falla el intervalo de confianza de la aproximación normal.
\textcite{brown2001interval} mostraron que su cobertura real oscila de forma
irregular, incluso con muestras grandes, y empeora con proporciones cercanas a 0
o a 1, y recomiendan el intervalo de \textcite{wilson1927probable}. Sea $n$ el
número de misiones, $\hat{p}$ la proporción observada de éxitos y $z$ el cuantil
de la confianza deseada ($z = 1{,}96$ para el 95~\%). El intervalo de Wilson es
\begin{equation}
  \frac{\hat{p} + \dfrac{z^{2}}{2n} \pm z\sqrt{\dfrac{\hat{p}(1-\hat{p})}{n} + \dfrac{z^{2}}{4n^{2}}}}{1 + \dfrac{z^{2}}{n}} .
\end{equation}
Su ancho decide qué se puede afirmar con una muestra dada.

El segundo principio es de método. Para \textcite{popper1959logic}, una
afirmación aporta conocimiento empírico solo si alguna observación posible podría
refutarla. Trasladado a un experimento, un criterio de aceptación solo es
evidencia si alguna ejecución del sistema puede incumplirlo.

La demostración con los agentes físicos sigue otra lógica. La norma ASTM
F3244-21 \autocite{astm2021f3244} establece un método de prueba para la
navegación de robots en un área definida, con repeticiones del recorrido entre
marcas fijas en el piso. El número de repeticiones fija cuánto se puede afirmar:
si las $n$ repeticiones salen bien, la fiabilidad mínima que se demuestra con una
confianza $C$ es
\begin{equation}
  R = (1-C)^{1/n} ,
\end{equation}
que con $n = 10$ y $C = 0{,}85$ da $R \approx 0{,}83$. Con las cinco a diez
repeticiones que caben en el calendario, la campaña física demuestra que el
sistema opera con los agentes físicos, pero no estima su tasa de éxito con la precisión de la simulación.

% =====================================================================
\section{Marco legal}
\label{sec:legal}
% =====================================================================

El proyecto se desarrolla dentro de la política pública colombiana de
transformación digital y de ciencia, tecnología e innovación, que enmarcan dos
documentos. El CONPES 3975 \autocite{conpes3975} fija la Política Nacional para la
Transformación Digital e Inteligencia Artificial. El CONPES 4069
\autocite{conpes4069} establece la Política Nacional de Ciencia, Tecnología e
Innovación. Ambos piden a la academia que desarrolle prototipos con herramientas
de la Cuarta Revolución Industrial, entre ellas la robótica autónoma, y que lo
haga con diseño local e impacto social.

Este trabajo encaja en esa línea. Es un prototipo de asistencia robótica
construido sobre plataformas de bajo costo y software abierto, dentro de las
limitaciones presupuestales de un entorno universitario.



\subsection{Alcance de la operación y responsabilidad}

El alcance del proyecto excluye de forma explícita que los agentes interactúen mecánicamente con el edificio. No accionan puertas ni botones de ascensor, no
transportan carga y no cruzan las zonas de transición vertical. Esa delimitación
se adoptó por razones técnicas y de costo, y tiene además una consecuencia
normativa favorable: evita intervenir sistemas sujetos a regulación específica de
seguridad en edificaciones.



% =====================================================================
\section{Marco ambiental}
\label{sec:ambiental}
% =====================================================================

El impacto ambiental del proyecto es mínimo. Los dos agentes son equipo institucional que ya existía: el proyecto no los compró, y el presupuesto del anteproyecto los contabiliza como uso y desgaste. Funcionan con baterías pequeñas y de bajo consumo, y una sesión de campo completa gasta una fracción de su carga. El entorno de evaluación es el propio edificio de la universidad, sin infraestructura dedicada ni modificaciones. El único cuidado que exige el prototipo es el de las baterías de polímero de litio de la tracción, que al final de su vida útil deben ir a los canales de disposición de residuos electrónicos de la institución y no a la basura doméstica.



% VERIFICAR: si la universidad tiene un programa institucional de gestión de RAEE
% (residuos de aparatos eléctricos y electrónicos), nombrarlo aquí en concreto.
% No se afirma que exista sin comprobarlo.
